% ---------------------------------------------------------------------
% Data availability statement + dataset description for the IEEE Access
% manuscript.  Place the statement immediately before the references;
% the description belongs either in the statement itself (if short) or
% as an appendix.
%
% PLACEHOLDERS to fill before submission -- marked \TODO:
%   - the DOI and repository name
%   - the archive size
%   - the licence
% ---------------------------------------------------------------------

\providecommand{\TODO}[1]{\textcolor{red}{[\textbf{TODO:} #1]}}

\section*{Data Availability}

The excitation dataset, the load-cell reference measurements and the
complete analysis code are openly available at
\TODO{repository and DOI}, released under \TODO{licence}.

The archive contains the $140$ gate-passing excitation runs of
Sec.~VIII-A, covering five centre-of-mass configurations on two axes in
both tip directions at four commanded ramp rates. Each run carries the
three sensing channels used here --- individual rotor speeds, from
which the applied moment and collective thrust are reconstructed; the
onboard rate gyro; and the motion-capture pose --- together with the
identified onset and the gate outcomes that produced it. The load-cell
offset measurement that serves as ground truth is included for all ten
case/axis configurations, as are the two-stage calibration constants
and the commanded moment offsets flown in Sec.~VIII-E.

Three uses are anticipated beyond reproducing the results reported
here, and the archive is organised for them.

\emph{As a benchmark for onset detection.} Every run is paired with an
independent measurement of the quantity the onset is used to estimate,
so a detector can be scored on the delivered offset and not only on
self-consistency. The seven detectors compared in
Table~\ref{tab:bench} are included with the harness that runs them, so
a new method can be added and evaluated against the same ten
configurations without re-deriving the inversion.

\emph{As a record of compliant ground contact.} The motion-capture arcs
resolve the pivot arm to a few millimetres and bound the rotation-centre
height and the contact migration during the tip, quantities that are
otherwise difficult to obtain for a landing gear under load. Sec.~VI-B
uses them to bound the ground-contact budget; they are reported in full
so that a different contact model can be fitted to them.

\emph{As a testbed for tip-over dynamics.} The ramp rate spans a factor
of twelve, which separates the rate-dependent and excursion-dependent
terms of any model of the departure, not only the one used here. The
raw traces are provided so that alternative dynamic models can be
identified from the same events.

One limitation should be stated with the release. The excitation was
performed at a single rotor clearance, so the aerodynamic and
contact-geometry contributions to the static threshold deficit are
degenerate in this dataset and cannot be separated by any analysis of
it; Sec.~VI-B reports three independent attempts and the reason each
fails. Separating them requires a clearance sweep, which the archive
does not contain.

\TODO{state the archive size and format, e.g. ``The archive is N GB;
raw logs are provided as rosbags with a Parquet mirror of the derived
per-run quantities.''}
